% TOP_PROBLEM: {"id": 1347, "title": "Partition Principle Implies Axiom of Choice", "category_id": 10, "category": {"id": 10, "name": "set_theory", "display_name": "Set Theory", "description": "Foundations of mathematics, infinite sets, and cardinality.", "slug": "set-theory", "order_index": 10, "created_at": "2026-07-31T15:26:25.671Z"}, "status": "open"}
% FIRST_POSTED: 2026-10-06
% ATTEMPT_STATUS: solved
% OPENAI_CLAIM: {"schema_version":1,"source_repo":"https://github.com/openai/math","source_commit":"adc7f1241b42e322a6451854ab7e4b4c146bf78a","release_date":"2026-10-06","match":"full","resolution":"disproved","independently_reviewed":false,"second_pass":"agreed","families":[{"family":"244","title":"The Partition Principle does not imply Choice","subject":"Mathematical logic","manuscripts":[{"title":"The Partition Principle does not imply Choice","pdf_path":"preprints/The-Partition-Principle-does-not-imply-Choice-September-24-2026/partition-principle-without-choice.pdf","readme_path":"preprints/The-Partition-Principle-does-not-imply-Choice-September-24-2026/README.md","date":"2026-09-24","theorem_ref":"thm:main-consistency (sections/introduction.tex:23); thm:main-transitive (sections/introduction.tex:29)","statement_sources":[{"path":"preprints/The-Partition-Principle-does-not-imply-Choice-September-24-2026/build/sections/introduction.tex","line":23,"label":"thm:main-consistency"},{"path":"preprints/The-Partition-Principle-does-not-imply-Choice-September-24-2026/build/sections/introduction.tex","line":29,"label":"thm:main-transitive"}]}],"lean":{"doc_path":"lean/docs/244.md","comparators":["lean/ComparatorChallenges/PartitionConsistency.json","lean/ComparatorChallenges/PartitionConsistency.lean","lean/ComparatorChallenges/PartitionPrinciple.json","lean/ComparatorChallenges/PartitionPrinciple.lean"],"declarations":["OAI.PartitionPilot.Forcing.TransitiveGround.exists_model_partitionPrinciple_without_choice"],"covers_main_theorem":true},"reasoning_trace":null}]}
\documentclass[11pt]{article}
\usepackage[margin=25mm]{geometry}
\usepackage{fontspec}
\setmainfont{Latin Modern Roman}
\usepackage{amsmath,amssymb,mathtools}
\usepackage{unicode-math}
\setmathfont{Latin Modern Math}
\usepackage{enumitem}
\usepackage{xurl,hyperref,bookmark}
\hypersetup{colorlinks=true,urlcolor=blue}
\setcounter{secnumdepth}{0}
\setlength{\parindent}{0pt}
\setlength{\parskip}{5pt}
\setlength{\emergencystretch}{3em}
\begin{document}
\section*{Partition Principle Implies Axiom of Choice}

\subsection{OpenAI's claim}
OpenAI presents the following claims relevant to this question, with manuscript numbers following paper order in Sources. In manuscript 1 (family 244; 2026-09-24), the asserted result is that relative to consistency of ZF, the Partition Principle can hold together with ordinal-indexed choice and failure of full Choice.

\subsection{Mathematical statement of the claimed result}
\paragraph{Manuscript 1.}
The Partition Principle $\mathrm{PP}$ means that a surjection $X\twoheadrightarrow Y$ implies some injection $Y\hookrightarrow X$, not necessarily a section. Let $\mathrm{AC}_{\rm wo}$ be choice for every ordinal-indexed family of nonempty sets. The claimed relative consistency is
\[
 \operatorname{Con}(\mathrm{ZF})\Rightarrow
 \operatorname{Con}(\mathrm{ZF}+\mathrm{PP}+\mathrm{AC}_{\rm wo}+\neg\mathrm{AC}).
\]
More concretely, from a countable transitive $V\models\mathrm{ZFC}$ the manuscript claims set forcing and a transitive symmetric intermediate model $V\subseteq W\subseteq V[G]$ with the same ordinals and $W\models\mathrm{ZF}+\mathrm{PP}+\mathrm{AC}_{\rm wo}+\neg\mathrm{AC}$. Every countable sequence of elements of $V$ that lies in $W$ already lies in $V$. Thus PP is claimed not to imply full Choice over ZF at this relative-consistency strength.

\subsection{Outline of the argument}
For manuscript 1 (family 244; 2026-09-24), a symmetric forcing construction creates nonwellorderable generic labels and orbit maps that present arbitrary sets by label--ordinal products. Controlled supports give ordinal-family choices, while quotient classification supplies injections opposite surjections without making the labels wellorderable.

\subsection{Scope relative to this problem}
The displayed question has a negative answer under the manuscript's stated construction. Its counterexample meets the conditions in this problem and defeats the universal assertion, so this is a full disproof claim. The supporting result is manuscript 1 (family 244; 2026-09-24).

\subsection{Formal verification}
For manuscript 1 (family 244; 2026-09-24), the documented Lean scope is the relative-consistency implication for ZF, the Partition Principle, ordinal-indexed choice, and failure of Choice; a separate transitive-model construction has an inaccessible-cardinal ground assumption and does not encode all the paper's preservation assertions.

The selected statement includes the clause used for this problem.

The linked family Comparator checks are \texttt{PartitionConsistency.lean}, \texttt{PartitionPrinciple.lean}.

Their named declarations include \texttt{OAI.PartitionConsistency.main\_consistency}, \texttt{OAI.PartitionPilot.Forcing.TransitiveGround.exists\_model\_partitionPrinciple\_without\_choice}.

This site has not executed the Lean code. Described statement coverage is not an independent proof check.

\subsection{Status caveat}
The manuscripts were produced by an internal OpenAI model. OpenAI cautions that results without formalization can contain mistakes. This site has not independently checked the proofs; the label records the scope of OpenAI's claim.

\subsection{Sources}
\begin{itemize}
\item OpenAI, \emph{The Partition Principle does not imply Choice}, 2026-09-24; citation key \texttt{OAI:The-Partition-Principle-does-not-imply-Choice-September-24-2026}. \href{https://github.com/openai/math/blob/adc7f1241b42e322a6451854ab7e4b4c146bf78a/preprints/The-Partition-Principle-does-not-imply-Choice-September-24-2026/partition-principle-without-choice.pdf}{PDF (GitHub)}; \href{https://github.com/openai/math/raw/adc7f1241b42e322a6451854ab7e4b4c146bf78a/preprints/The-Partition-Principle-does-not-imply-Choice-September-24-2026/partition-principle-without-choice.pdf}{PDF (direct)}; \href{https://github.com/openai/math/blob/adc7f1241b42e322a6451854ab7e4b4c146bf78a/preprints/The-Partition-Principle-does-not-imply-Choice-September-24-2026/README.md}{Manuscript page and citation}; \href{https://github.com/openai/math/blob/main/preprints/The-Partition-Principle-does-not-imply-Choice-September-24-2026/partition-principle-without-choice.pdf}{Latest version}; Statement source: \href{https://github.com/openai/math/blob/adc7f1241b42e322a6451854ab7e4b4c146bf78a/preprints/The-Partition-Principle-does-not-imply-Choice-September-24-2026/build/sections/introduction.tex\#L23}{thm:main-consistency}; Statement source: \href{https://github.com/openai/math/blob/adc7f1241b42e322a6451854ab7e4b4c146bf78a/preprints/The-Partition-Principle-does-not-imply-Choice-September-24-2026/build/sections/introduction.tex\#L29}{thm:main-transitive}.
\item Family 244: \href{https://github.com/openai/math/blob/adc7f1241b42e322a6451854ab7e4b4c146bf78a/lean/docs/244.md}{Lean scope}.
\item \href{https://github.com/openai/math/blob/adc7f1241b42e322a6451854ab7e4b4c146bf78a/lean/ComparatorChallenges/PartitionConsistency.json}{Comparator PartitionConsistency.json}.
\item \href{https://github.com/openai/math/blob/adc7f1241b42e322a6451854ab7e4b4c146bf78a/lean/ComparatorChallenges/PartitionConsistency.lean}{Comparator PartitionConsistency.lean}.
\item \href{https://github.com/openai/math/blob/adc7f1241b42e322a6451854ab7e4b4c146bf78a/lean/ComparatorChallenges/PartitionPrinciple.json}{Comparator PartitionPrinciple.json}.
\item \href{https://github.com/openai/math/blob/adc7f1241b42e322a6451854ab7e4b4c146bf78a/lean/ComparatorChallenges/PartitionPrinciple.lean}{Comparator PartitionPrinciple.lean}.
\item \href{https://github.com/openai/math/blob/adc7f1241b42e322a6451854ab7e4b4c146bf78a/CONTENTS.md}{OpenAI catalogue entry}.
\end{itemize}
\end{document}
